\documentclass[tikz,border=5pt]{standalone}
\usetikzlibrary{spath3,intersections}
\begin{document}
    \tikzset{
        bridging path/.initial=arc, % 定义键的初始值 path=`/tikz/`
        bridging span/.initial=8pt, 
        bridging gap/.initial=4pt, 
        bridge/.style 2 args={
            % #1: over path #2: under path
            spath/split at intersections with={#1}{#2}, 
            spath/insert gaps after components={#1}{\pgfkeysvalueof{/tikz/bridging span}}, 
            spath/join components upright with={#1}{\pgfkeysvalueof{/tikz/bridging path}}, 
            spath/split at intersections with={#2}{#1}, 
            spath/insert gaps after components={#2}{\pgfkeysvalueof{/tikz/bridging gap}},
        }
    }
    % If used in the preamble, this needs surrounding in \AtBeginDocument
    %\AtBeginDocument{%
        \tikz[overlay] { % 保存全局路径
            \path[spath/save global=arc] 
                (0,0) arc[radius=1cm,start angle=180, delta angle=-180]; 
        }%
    %}
    \begin{tikzpicture}
            \path[spath/save=over] (0,0) -| ++(1,1) -| ++(-1,1) -| ++(1,1) -| ++(-1,1);
            \path[spath/save=under] (.5,-.5) -- ++(0,4);
            \tikzset{bridge={over}{under}} % 批量控制两个路径的 arc 路径
            \draw[spath/use=over]; 
            \draw[spath/use=under]; 
    \end{tikzpicture}
\end{document}